% Expose IX, 1.0--1.1.
% French transcription lines 144--260; authority scan physical pp. 331--334
% (source folios 319--322).  Ordinary pages were checked in four overlapping
% 1100-dpi bands.  The missing prime on A'_o in the final sentence of 1.1 is
% an ordinary printed/source slip and is restored below.

\medskip
\noindent
\underline{1. Neron models of abelian schemes: notation; the canonical pairing}
\[
 \Phi_o\times\Phi'_o\longrightarrow(\mathbb Q/\mathbb Z)_k,
\]
\underline{and the canonical biextension $W^o$ of $(A^o,{A'}^o)$.}

\medskip
\noindent
1.0. Throughout this expos\'e we fix a trait $S$, for which we use the
usual notation $s$, $\eta$, $k$, $K$, and so forth, of I~1.  Let $A_K$
be an abelian scheme over $K$, and denote its dual abelian scheme [5] by
$A'_K$.  The duality between $A_K$ and $A'_K$ is given by a divisorial
correspondence $W_K$ on $(A_K,A'_K)$, that is, an invertible sheaf on
$A_K\times A'_K$, birigidified with respect to the identity sections of
$A_K$ and $A'_K$.  It is determined up to unique isomorphism by its
class (cf. VIII~3.1)
\begin{equation}\tag{1.0.1}
 w_K\in\mathrm{Biext}^1(A_K,A'_K;\mathbb G_{mK}).
\end{equation}
An important role will be played by the corresponding pairing
(VIII~2.2)
\begin{equation}\tag{1.0.2}
 \varphi:T_\ell(A_K)\times T_\ell(A'_K)
 \longrightarrow\mathbb Z_\ell(1)_K
 \overset{\mathrm{def}}{=}T_\ell(\mathbb G_{mK}),
\end{equation}
where $\ell$ is an arbitrary prime number.  When
$\ell\ne p=\operatorname{char}k$, giving $\varphi$ is equivalent to
giving a pairing of finite free $\mathbb Z_\ell$-modules
\begin{equation}\tag{1.0.3}
 \varphi:T_\ell(A_K(\overline K))\times
 T_\ell(A'_K(\overline K))\longrightarrow T_\ell(\overline K^*),
\end{equation}
compatible with the action of
\[
 \pi=\operatorname{Gal}(\overline K/K)
\]
on the three modules in question, where $\overline K$ is a separable
closure of $K$.  It is known (VIII~3.2) that this pairing is a perfect
duality: it identifies each of the modules
$T_\ell(A_K(\overline K))$ and $T_\ell(A'_K(\overline K))$ with the
dual of the other, with values in the invertible module
$\mathbb Z_\ell(1)=T_\ell(\overline K^*)$.

At times we shall suppose that a polarization of $A_K$ has been chosen;
thus we are given an element
\begin{equation}\tag{1.0.4}
 \xi\in\operatorname{NS}(A_K/K)
 =\underline{\operatorname{NS}}_{A_K/K}(K),
\end{equation}
where
\[
 \underline{\operatorname{NS}}_{A_K/K}
 \overset{\mathrm{def}}{=}
 \underline{\operatorname{Pic}}_{A_K/K}/
 \bigl(\underline{\operatorname{Pic}}{}^o_{A_K/K}=A'_K\bigr).
\]
In the usual way it defines a homomorphism
\begin{equation}\tag{1.0.5}
 \psi_\xi:A_K\longrightarrow A'_K,
\end{equation}
hence a homomorphism
\begin{equation}\tag{1.0.6}
 T_\ell(\psi_\xi),\quad\text{or simply}\quad
 \psi_\xi:T_\ell(A_K)\longrightarrow T_\ell(A'_K),
\end{equation}
and a bilinear form $\varphi_\xi$
\begin{equation}\tag{1.0.7}
 \begin{split}
 \varphi_\xi(x,y)&=\varphi(x,\psi_\xi(y)),\\
 \varphi_\xi:T_\ell(A_K)\times T_\ell(A_K)
 &\longrightarrow\mathbb Z_\ell(1)_K,
 \end{split}
\end{equation}
which, for $\ell\ne p$, may also be interpreted as a pairing of
finite free $\mathbb Z_\ell$-modules with operators:
\begin{equation}\tag{1.0.8}
 \varphi_\xi:T_\ell(A_K(\overline K))\times
 T_\ell(A_K(\overline K))\longrightarrow T_\ell(\overline K^*).
\end{equation}
These constructions make sense for every element of the
Neron--Severi group of $A_K$, and yield an \underline{alternating} form
$\varphi_\xi$, in either interpretation (1.0.7) or (1.0.8)
(VIII~2.2.11).  The fact that $\xi$ is a \underline{polarization}
implies that $\psi_\xi$ in (1.0.5) is an isogeny; equivalently, since
$\dim A_K=\dim A'_K$, it is an epimorphism of groups.  This is also
equivalent to the injectivity of $T_\ell(\psi_\xi)$ in (1.0.6) (or to
its having finite cokernel), and finally to the form $\varphi_\xi$ in
(1.0.7) being \underline{separating}.

\medskip
\noindent
1.1. Recall that there exists a \underline{smooth scheme} $A$ over
$S$, called the \underline{Neron model of $A_K$}, which represents the
functor
\begin{equation}\tag{1.1.1}
 S'\longmapsto\operatorname{Hom}(S'_\eta,A_K)
\end{equation}
on the category of smooth $S$-schemes.  Thus, by definition, there is
an isomorphism, functorial in the smooth relative $S$-scheme $S'$,
\begin{equation}\tag{1.1.2}
 \operatorname{Hom}_S(S',A)\simeq
 \operatorname{Hom}_K(S'_\eta,A_K).
\end{equation}
Neron [19] constructed $A$ when the residue field $k$ of $S$ is
perfect.  Raynaud extended the construction to the general case in an
unwritten seminar held at IHES in 1966--67 (see [21] for the statement
of Raynaud's principal results, which complete those of Neron in
loc. cit.).  We shall admit this existence result here, together with
the fact that $A$ is a quasi-projective \underline{scheme} over $S$.
Applying (1.1.2) when $S'$ is smooth \underline{over} $\eta$, one also
finds that the following canonical morphism is an isomorphism:
\begin{equation}\tag{1.1.3}
 A\times_S\operatorname{Spec}(K)\xrightarrow{\ \sim\ }A_K.
\end{equation}
This justifies the paired notation $A,A_K$; when convenient we shall
also write $A_\eta$ in place of $A_K$.  Using (1.1.3) as an identifying
isomorphism, the canonical bijection (1.1.2) is simply the map given by
\underline{restriction to the generic fibers}.

Since the functor (1.1.1) plainly has the structure of a commutative
group, the Neron model $A$ representing it itself has the structure of
a \underline{smooth commutative group scheme over $S$}.  The remainder
of this expos\'e is devoted to the study of this group scheme.

An important invariant arising from the Neron model is, of course, the
smooth commutative group scheme of finite type over $k$
\begin{equation}\tag{1.1.4}
 A_o=A\times_Ss=A_k.
\end{equation}
Following the general notation, we denote its identity component by
$A_o^o$ and also consider the group scheme
\begin{equation}\tag{1.1.5}
 \Phi_o=A_o/A_o^o
\end{equation}
of connected components of $A_o$, which is finite and etale over $k$.
In terms of a separable closure $\overline k$ of $k$, it is therefore
determined by the ordinary finite group
\begin{equation}\tag{1.1.6}
 \Phi_o(\overline k),
\end{equation}
together with the natural action of
\[
 \pi_o=\operatorname{Gal}(\overline k/k)
\]
on this group.

Along with the Neron model $A$ of $A_K$, we shall study
\begin{equation}\tag{1.1.7}
 A'=\text{the Neron model of }A'_K,
\end{equation}
which is likewise a smooth commutative group scheme of finite type over
$S$.  We shall denote by
\begin{equation}\tag{1.1.8}
 \Phi'_o=A'_o/{A'}_o^o
\end{equation}
the group of connected components of its special fiber $A'_o$.
